\subsection{Transitivity of algebraic extensions. Fields that are relatively algebraically closed in an extension field}

PROPOSITION 3. --- Let E and F be two extension fields of a field K such that \(K \subset E \subset F\) . For F to be algebraic over K it is necessary and sufficient that E should be algebraic over K and F algebraic over E.

The condition is necessary, by V, p.~17, Cor. 2. Let us show that it is sufficient. Let x be any element of F; it is algebraic over E; let \(g \in E[X]\) be its minimal polynomial over E. If A is the (finite) set of coefficients of g, then \(g \in K(A)[X]\) , hence x is algebraic over \(K(A)\) and \(K(A \cup \{x\}) = K(A)(x)\) is of finite degree over \(K(A)\) . Now \(A \subset E\) and E is algebraic over K, hence \(K(A)\) is of finite degree over K, by Th. 2. Therefore (V, p.~10, Th. 1), \(K(A \cup \{x\})\) is of finite degree over K, which shows x to be algebraic over K (V, p.~18, Prop. 2).

DEFINITION 2. --- A subfield K of a field E is said to be relatively algebraically closed in E if every element of E which is algebraic over K, belongs to K.

It comes to the same to say that K is the only algebraic extension of K contained in E. Every field K is relatively algebraically closed in itself. In § 4 we shall study the fields which are relatively algebraically closed in each extension field.

PROPOSITION 4. --- Let E be an extension of a field K; the set L of elements of E which are algebraic over K forms a subextension of E which is relatively algebraically closed in E.

For the field \(\mathbf{K}(\mathbf{L})\) is algebraic over \(\mathbf{K}\) (Cor. 1 of Th. 2), hence \(\mathbf{K}(\mathbf{L}) \subset \mathbf{L}\) ; it follows that \(\mathbf{K}(\mathbf{L}) = \mathbf{L}\) and \(\mathbf{L}\) is a field. On the other hand, if \(x \in \mathbf{E}\) is algebraic over \(\mathbf{L}\) , it is so also over \(\mathbf{K}\) (Prop. 3) and hence belongs to \(\mathbf{L}\) .

The extension L of K consisting of all the elements of E that are algebraic over K is called the relative algebraic closure of K in E; it is the largest algebraic extension of K contained in E.

